\section{Exploitation Playbooks}\label{sec:playbooks}
% <EDIT> One playbook per vulnerability class present in the codebase. Each documents a real,
% reachable sink confirmed in the code, with a concrete (illustrative) PoC, the impact, and the
% mitigation. PoCs are for authorised testing only.
\subsection{\fillme{Vulnerability Class, e.g. SQL Injection}}
\textbf{Severity: \critical\quad Auth: \fillme{None (incl. login-path)}}

\textbf{Description.} \fillme{One sentence describing the sink and how it is reached, referencing
the systemic root cause it stems from, e.g. "(Root cause R1)".}
\begin{codebox}[\fillme{PoC -- e.g. auth bypass / extraction via injectable parameter}]
# illustrative request/command only -- replace with your own, no real hostnames/tokens/credentials
\end{codebox}
\noindent\textbf{Representative sink:} \fillme{module/path/File.java:NN}\quad\textbf{Endpoint:} \fillme{METHOD /path}

\textbf{Impact.} \fillme{One sentence.} \textbf{Mitigation.} \fillme{One sentence.}

\subsection{\fillme{Vulnerability Class 2}}
\textbf{Severity: \fillme{\high}\quad Auth: \fillme{None}}

\textbf{Description.} \fillme{...}
\begin{codebox}[\fillme{PoC title}]
... illustrative PoC command/request goes here ...
\end{codebox}
\noindent\textbf{Representative sink:} \fillme{module/path/File.java:NN}\quad\textbf{Endpoint:} \fillme{METHOD /path}

\textbf{Impact.} \fillme{...} \textbf{Mitigation.} \fillme{...}
